\chapter{Неуспешни експерименти, безбедност и ограничувања}
\label{chap:failures}

\section{Неуспехот како резултат}

Во роботика е лесно да се прикаже успешниот command path и да се изостави
отсуството на ефект. Во овој проект неуспехот е evidence ако е датиран,
ограничен и корелиран. \Cref{tab:failures} ги сумира клучните случаи; подолу
се анализираат причините и преостанатата неизвесност.

\begin{table}[htbp]
\centering
\caption{Избрани неуспешни или делумни експерименти}
\label{tab:failures}
\small
\begin{tabularx}{\textwidth}{L{0.22\textwidth}YY}
\toprule
\textbf{Случај} & \textbf{Набљудуван резултат} & \textbf{Научена лекција/статус} \\
\midrule
Три ROS height trials & Nonzero intent и делумна wire evidence, но нулто операторски забележано движење & Publication не е actuation; не се зголемува амплитуда без да се затвори ownership/acceptance границата. \\
Auto-mode neutral height & $-156..+260$ packets, joint span 0.00084 rad, без значајно дишење & Mode switch не го направил малиот profile видлив. \\
Periodic Retroid height & Whole-run span вклучувал mode transition; post-handshake periodic span само 0.0011444 rad & Measurement window мора да ја изолира активната фаза. \\
Current-pose SDK takeover & Backward stepping и големи joint spans од микроскопски offsets & \code{RobotStateInit} не е current-pose acquisition; hard gate останува. \\
Stale telemetry & Stand или expression запрел на old 100 ms rule & Двостепен pause/recovery 100/250 ms со 20 fresh frames. \\
Joy/Sadness/Fear unload miss & Target paw задржал премногу load & Gate не се заобиколува; landing/recovery се довршуваат. \\
Sadness deep bow & Pitch 10.006/10.005 degrees & Attitude gate правилно го прекинал unsafe candidate. \\
Anger repeat & Третиот cycle relatch-ирал само три supports & Canonical support reset, без слабеење на gate. \\
Battery floor & 75\% floor видел 74\% & Immediate release; подоцнежен 25\% run само со operator decision. \\
OpenAI/DNS timeout & APIConnectionError и client timeout пред provider window & Explicit resolver и deadline derived from retry window; robot safety path останал независен. \\
Gazebo clock stall & Longer idle sweep запрел во neutral simulated time & Runtime failure е евидентиран; не се претставува како animation pass. \\
Visual mismatch & Безбедни Joy/Sadness/Fear прототипи не изгледале како емоцијата & Telemetry pass не е visual acceptance. \\
\bottomrule
\end{tabularx}
\end{table}


\section{ROS команда без физичко движење}

Очекувањето било height-only Joy или Neutral waveform да го помести трупот.
Забележани биле emotion state, mapper output и \code{/simple\_cmd}; по bridge
serialization fix се виделе 39 height datagrams, 10 heartbeats и final exact
zero на wire. Сепак joint span бил околу 0.00099 rad и немало визуелно
движење. Потврдена причина за дел од packet loss била queue depth 1 и
back-to-back uncommissioned axes, но \code{jy\_exe} acceptance на nonzero
height останал нерешен. Fail-safe исходот бил exact zero, DISARMED, restored
transmit flags false и без повторно зголемување на амплитуда.

Direct Retroid-compatible test подоцна докажал дека exact mode ordering,
cadence и companion value се важни. Но дури и таму првиот whole-run joint span
бил погрешно толкуван додека не се одделил Move-to-Pose transition. Ова ја
нагласува потребата baseline/window да се дефинира пред активната команда.

\section{Небезбеден takeover}

Очекувањето било MotionSDK да ја преземе моменталната стоечка поза и да додаде
мал offset. Набљудувани биле backward stepping и големи joint excursions.
Source/manual review потврдил дека \code{RobotStateInit} reset-ира кон zero;
„current-pose acquisition“ била погрешна хипотеза. Test е запрен пред breathing,
takeover hard gate не е bypass-иран, а официјалниот пат е пренасочен кон
state-1 zero-to-stand sequence под соодветна restraint. Овој негативен
резултат е една од најважните безбедносни лекции во проектот.

\section{Stale feedback и watchdog recovery}

Еден early official-breathing обид завршил со
\code{ABORT stale 0x0906 feedback}. Причината не била robot-state fault, туку
пасивни scheduling gaps околу 100--166 ms. Наместо watchdog да се исклучи,
воведено е двостепено однесување: trajectory-time freeze над 100 ms, hold на
последна валидна поза и resume по 20 fresh frames; на 250 ms се release-ира.
Во подоцнежните тестови pauses се забележани и опоравени, додека state 8
сепак веднаш активирал release. Значи толеранцијата кон bounded jitter не ја
претвора мртвата телеметрија во дозвола за blind control.

\section{Contact unload, landing и support miss}

Joy, Sadness, Fear и Anger прототипите покажале дека joint trajectory не е
доказ дека стапалото е unload-ирано. Примери се 8.875 N Sadness, 7.264 N Fear и
12.522 N повторен Joy left paw. Во секој случај paw бил спуштен и support бил
повторно оценет. Threshold не бил намалуван за да се добие pass.

Кај Anger, final right landing од третиот repeat имал target load 13.133 N,
но само три support latches. Motion визуелно изгледал добро, сепак endurance
резултатот останал failure. X/Y recenter не помогнал; дури exact canonical
support reset меѓу paws поминал шест placements. Ова покажува дека force value
на target leg и целокупната support topology мора да се разгледуваат заедно.

\section{Battery и attitude interlocks}

Sadness кандидати од 70 mm и 56 mm стигнале малку над $10^{\circ}$ pitch и
биле прекинати пред анимираната фаза. Вредностите 10.006 и 10.005 степени не се
„заокружуваат“ во pass: hard bound е бинарен. Final 48 mm front-only design го
вратил margin-от без да се менува gate.

Во Anger-to-Disgust сесија explicit 75\% floor видел 74\% и веднаш release-ирал.
Retarget не бил броен како pass. Подоцнежна independent сесија користела веќе
конфигуриран 25\% project floor само по експлицитна operator одлука; другите
gates не биле ослабени. Овој пример покажува дека резултатот мора да го наведе
точниот floor, а не само „battery was healthy“.

\section{Network, chat и simulator failures}

Live OpenAI physical attempt не бил pass кога user appraisal стигнал како Joy,
но client timeout-ирал пред assistant response. Изолираната причина била
intermittent DNS resolution failure. Поправката и два live smoke tests се
извршени пред повторна физичка сесија. Во Gazebo, две longer animation sweep
извршувања запреле поради stalled simulated clock. Пократкиот all-category
sweep поминал, но тоа не го претвора stalled run во успех.

\section{Безбедносни и етички аспекти}

Безбедноста е системско својство, не придавка. Во овој проект се применуваат:

\begin{itemize}
  \item explicit current-task authorization, fresh telemetry и human operator;
  \item clear area, manufacturer restraint/hoist каде што е барано, независен
  STOP и без луѓе во опасната зона;
  \item еден exclusive sender и забрана за concurrent Retroid/SDK paths;
  \item state, error, battery, attitude, tracking, feedback и contact gates;
  \item deadline-bounded quintic trajectories, zero torque feed-forward и
  deterministic release;
  \item no-secret practice за API key, SSH и packet captures; и
  \item искрено известување за negative и unvalidated results.
\end{itemize}

Емоционалното движење може да го засили anthropomorphism. Затоа трудот не
тврди дека роботот „чувствува“ емоции; системот пресметува и изразува
категоризирана affective state. Корисникот треба да може да разликува engine
state, physical fallback и safety override. Status contract-от токму затоа не
ја преименува unsupported emotion во Neutral.

\section{Ограничувања}

Проектот е prototype на еден Lite3 и едно проверено software/network
опкружување. Нема формална human-subject study за препознатливоста; operator
acceptance е инженерски acceptance criterion, не популациски резултат. Contact
estimator-от е изведен од torque и model Jacobian, а не од независни force
sensors. Battery floors се разликувале по authorized session. MotionSDK
takeover и release може да завршат со vendor sit/lie behavior; не е seamless
handoff.

Физичката allowlist е Neutral-only. Четири профили се планирани, а persistent
Joy/Anger contact-miss recovery сè уште нема live revalidation. Нема airborne
hop, locomotion, yaw, forceful stomp или general obstacle-aware behavior.
Gazebo Classic има завршен upstream life cycle, што е дополнителна долгорочна
техничка обврска. Овие ограничувања го спречуваат
проектот да се нарече завршена production platform.
